\chapter*{Introduzione}
\addcontentsline{toc}{chapter}{Introduzione}
\markboth{INTRODUZIONE}{INTRODUZIONE}

L'Agenzia Italiana del Farmaco (AIFA) regola la rimborsabilità di un sottoinsieme rilevante di farmaci attraverso le cosiddette \emph{Note AIFA}, prescrizioni testuali che fissano criteri di inclusione, esclusione ed eccezione subordinati alle caratteristiche cliniche del paziente. Una sola nota può coinvolgere indici di rischio cardiovascolare, livelli sierici di vari analiti, presenza simultanea di più diagnosi e cronologie terapeutiche, e la decisione finale (se il farmaco è a carico del Servizio Sanitario Nazionale o del cittadino) dipende dalla congiunzione di tutte queste condizioni. La verifica avviene tipicamente in fase di prescrizione, quando il medico ha tempo limitato e deve consultare il testo della nota in parallelo all'anamnesi. È un'attività a basso valore aggiunto cognitivo ma ad alto costo di attenzione: un errore in più o in meno si traduce in farmaco non rimborsato a carico del paziente, in segnalazioni amministrative al medico o, nei casi peggiori, in un danno clinico quando una controindicazione viene trascurata.

I sistemi di supporto alla decisione clinica (\emph{Clinical Decision Support Systems}, CDSS) nascono proprio per ridurre questo carico. Nella loro forma più diffusa codificano le regole in un motore deterministico e presentano al medico l'esito di una verifica binaria: è lo stato dell'arte per il controllo delle Note AIFA nei principali software gestionali italiani, ed è efficace finché le regole sono stabili e in numero contenuto. Restano però due limiti che pesano proprio nel momento della prescrizione. Il primo è la spiegazione: un esito \emph{non rimborsabile} restituito senza motivazione costringe il medico, di fronte a una decisione che giudica controintuitiva, a risalire da solo alla condizione che ha bloccato la prescrizione, riaprendo il problema cognitivo che il sistema dichiarava di risolvere. Il secondo è la verificabilità: chi deve controllare il sistema (l'azienda sanitaria, l'ente regolatore, lo stesso autore) ha bisogno di ispezionare l'intera catena che porta dalla regola al testo della nota, e un esito binario non offre questa tracciabilità.

I \emph{Large Language Model} (LLM) degli ultimi anni sembrano lo strumento naturale per colmare entrambi i limiti, soprattutto se combinati con tecniche di \emph{Retrieval-Augmented Generation} (RAG): un modello generativo può articolare in linguaggio naturale il motivo di una decisione e citare i passaggi della nota da cui quel motivo deriva. Affidare però a un LLM la decisione di rimborsabilità sarebbe imprudente sul piano clinico. Le garanzie di sicurezza richieste a un sistema prescrittivo non sono compatibili con un componente generativo: per quanto bassa sia la probabilità di errore su un dominio chiuso e documentato come quello delle Note AIFA, il suo output non è verificabile a priori contro una specifica scritta.

La tesi propone una via intermedia, di impianto neuro-simbolico, fondata su una separazione netta di responsabilità. La decisione di rimborsabilità è affidata in modo esclusivo a un componente simbolico deterministico, che traduce le Note AIFA in regole esplicite e ispezionabili: è qui, e solo qui, che risiedono le proprietà di sicurezza del sistema. Al modello generativo è riservato il solo compito di spiegare: riceve la decisione già presa e i passaggi delle Note pertinenti, e li riformula in un testo clinicamente leggibile, vincolato a citare le fonti che gli sono state fornite. La sicurezza della decisione è così garantita per costruzione; la qualità della spiegazione diventa invece una questione empirica, ed è la sua misura rigorosa il cuore del lavoro.

Il sistema è costruito e valutato su quattro Note AIFA (\notezero\ per gli inibitori di pompa protonica nella gastroprotezione cronica, \notetredici\ per le statine nella dislipidemia, \notesessanta\ per gli antinfiammatori non steroidei e \notenova\ per gli anticoagulanti orali diretti nella fibrillazione atriale non valvolare) scelte perché coprono insieme l'intera varietà di condizioni che una nota può imporre: soglie su valori continui, diagnosi multiple, eccezioni e interazioni fra parametri. La verifica empirica si fonda su un dataset \emph{gold} di \goldn\ casi clinici annotati manualmente, rispetto ai quali il comportamento del sistema viene confrontato con quello di due alternative naturali: un LLM lasciato a sé stesso e un LLM dotato di retrieval ma privo del componente simbolico.

Il valore scientifico del lavoro non sta tuttavia nell'architettura in sé, ma nelle domande a cui essa permette di rispondere. L'intera valutazione è organizzata attorno a sette domande di ricerca, d'ora in avanti abbreviate RQ (dall'inglese \emph{Research Question}) e numerate da RQ1 a RQ7, che fanno da filo conduttore a tutti i capitoli che seguono; ciascuna è formulata in modo da ammettere una risposta binaria sul comportamento osservato del sistema. Le prime due riguardano i due componenti presi singolarmente: il componente simbolico decide come l'annotatore umano sui casi gold (RQ1)? E il retrieval individua, per ogni caso, i passaggi delle Note che sostengono la decisione (RQ2)? La terza e la sesta riguardano la spiegazione: il testo generato dal modello si appoggia effettivamente alle fonti recuperate, senza introdurre informazioni esterne (RQ3)? E consente al medico di ricondurre ogni affermazione al testo della nota (RQ6)? La quarta interroga la robustezza del sistema rispetto a ripetizioni e a perturbazioni controllate dell'input (RQ4); la quinta ne misura il vantaggio rispetto alle due baseline (RQ5); la settima, infine, sintetizza in un giudizio complessivo l'utilità clinica e la conformità ai requisiti di sicurezza (RQ7). Rispondere a queste sette domande, con le riserve metodologiche che ciascuna porta con sé, è l'obiettivo scientifico del lavoro: tutto ciò che le precede, dal contesto clinico all'architettura del sistema, ne costituisce la premessa necessaria, il contorno che dà sostanza alle domande senza sostituirle.

Il contributo è triplice. Sul piano architetturale, la tesi mostra una strategia di integrazione neuro-simbolica praticabile per un CDSS prescrittivo, con un confine netto fra ciò che decide e ciò che spiega. Sul piano metodologico, propone un protocollo di valutazione strutturato per domande di ricerca, in cui ogni metrica è al servizio di una sola domanda binaria. Sul piano operativo, infine, rilascia un sistema riproducibile end-to-end, dal documento ufficiale AIFA fino ai numeri finali della valutazione.

Coerentemente con questa gerarchia, la tesi è organizzata in due parti. La prima, le \emph{premesse}, fornisce ciò che serve a porre le domande senza ancora rispondervi: il \Cref{ch:contesto} ricostruisce il contesto clinico e regolatorio delle Note AIFA e descrive le quattro Note oggetto di studio; il \Cref{ch:background} introduce gli strumenti tecnici (rule engine, pipeline RAG, approcci neuro-simbolici) e le metriche di valutazione che torneranno nei risultati; il \Cref{ch:progettazione} presenta l'architettura, motivando il principio di separazione fra decisione e spiegazione; il \Cref{ch:implementazione} descrive le scelte implementative concrete, dallo stack tecnologico ai principali algoritmi del rule engine e della pipeline di ingestione. La seconda parte è il cuore scientifico del lavoro: il \Cref{ch:impianto} fissa l'impianto della valutazione, ossia il gold standard e il protocollo sperimentale, ed enuncia le sette domande di ricerca; il \Cref{ch:valutazione} le affronta una per una, ciascuna in una sezione propria, con i risultati numerici e la discussione critica delle limitazioni residue; il \Cref{ch:conclusioni} ne sintetizza le sette risposte e indica gli sviluppi futuri. L'\Cref{app:riproducibilita} raccoglie i comandi e le versioni che permettono di rieseguire l'intera valutazione da PDF a numeri canonici.
